\chapter{ОСНОВНАЯ ЧАСТЬ}

\section{Проблематика разработки шахматного веб-приложения}

Разработка шахматного веб-приложения отличается от разработки обычного информационного сайта. В информационном сайте основная задача состоит в отображении материалов и обработке простых форм. В ChessView требуется синхронная работа двух игроков, проверка правил шахмат, хранение истории ходов, учет времени, восстановление соединения, обработка завершения партии и защита состояния от некорректных действий пользователя. Поэтому приложение должно проектироваться как полноценная распределенная система, а не как набор страниц.

Первая проблема связана с авторитетностью состояния. Браузер пользователя может быстро отрисовать ход на доске, но окончательное решение о допустимости хода должно принимать сервер. В ChessView backend получает UCI-ход, проверяет очередь хода, состояние часов и легальность действия через библиотеку python-chess, затем сохраняет новую FEN-позицию и только после этого отправляет событие \code{game_state}. Такой подход исключает ситуацию, при которой один клиент считает ход допустимым, а другой клиент получает другое состояние партии.

Вторая проблема связана с real-time обменом. REST API удобно использовать для профиля, списка турниров, истории партий и задач, однако live-партия требует двустороннего канала. Для этого в ChessView применяется WebSocket. Один endpoint \code{/ws?token=<access_token>} принимает события очереди, хода, предложения ничьей, чата и WebRTC signaling. Внутри backend событие разбирается по типу и передается доменному обработчику. Это снижает связанность кода и делает сетевой контракт понятным.

Третья проблема связана с хранением данных. База данных должна поддерживать не только пользователей и партии, но и зависимые сущности: ходы, сообщения, задачи, попытки решения, турниры, участников, раунды, пары, запланированные матчи, платежные события, административный аудит и проверки личности. Если хранить все в одной таблице или в неструктурированном JSON, дальнейшее развитие проекта станет сложным. Поэтому используется реляционная модель PostgreSQL с внешними ключами.

Четвертая проблема связана с границами ответственности между frontend и backend. Клиент должен быть удобным, быстрым и интерактивным, но не должен дублировать всю бизнес-логику сервера. В ChessView клиент выполняет маршрутизацию, отображает доску, отправляет действия, запускает локальный анализ Stockfish и показывает данные пользователя. Backend отвечает за авторизацию, партии, турниры, задачи, платежный эмулятор, административные действия и сохранение данных.

Пятая проблема связана с демонстрационной полнотой. Дипломный проект нельзя считать завершенным только по факту запуска главной страницы. В пояснительной записке нужно показать, какие пользовательские действия доступны, какие таблицы участвуют в этих действиях, какие API вызываются и какие ограничения есть у single-instance запуска. Поэтому структура работы раскрывает не только внешний вид, но и внутреннюю организацию приложения.



Дополнительно проблематику можно разложить по ролям системы. Для игрока важны быстрый вход, понятная доска, корректные часы и отсутствие спорных ситуаций по ходам. Для backend важны проверка токена, проверка права игрока ходить в конкретной партии, сохранение результата и рассылка события только участникам нужной комнаты. Для базы данных важны связи между users, games, moves и ratings-логикой. Для администратора важны управляемость аккаунтов, журнал действий и возможность видеть спорные записи. Такое разделение ролей использовано дальше при описании структуры ПО, БД, API, клиента и административной панели.

В ChessView нельзя переносить критичные решения на frontend. Клиентская часть работает в недоверенной среде: пользователь может изменить JavaScript в браузере, подменить запрос или отправить WebSocket-событие вручную. Поэтому backend должен повторно проверять все важные условия. Для хода это означает проверку game_id, user_id, стороны игрока, статуса партии, очереди хода, формата UCI и легальности хода. Для административного действия это означает проверку access token и роли admin. Для платежного сценария это означает проверку статуса payment intent и связанного пользователя.

Отдельная задача связана с понятностью архитектуры. Если смешать авторизацию, игру, турниры, задачи и админку в одном файле, проект станет трудно объяснять и сопровождать. Поэтому backend разделен на предметные домены, а frontend разделен на страницы, виджеты, features и shared-слой. В результате каждая часть проекта имеет собственную ответственность: game отвечает за игру, tournaments за турниры, puzzles за задачи, admin за административные операции, shared за общие инфраструктурные элементы.

Еще одна практическая проблема связана с отчетностью дипломного проекта. Руководитель и комиссия должны видеть не только итоговый интерфейс, но и внутреннее устройство: какие таблицы есть в БД, какие поля в них хранятся, какие API вызываются, как пользователь проходит сценарий, где расположен backend-код и как запускается приложение. Поэтому в пояснительной записке структура описывается раздельно: отдельно ПО, отдельно БД, отдельно API, отдельно клиент, отдельно админ-панель.

\section{Обзор существующих программных решений}

При проектировании ChessView были рассмотрены существующие шахматные платформы. Они показывают разные подходы к одной предметной области: массовая игровая платформа, бесплатный open-source сервис, турнирная система и учебная реализация. Сравнение приведено в таблице~\ref{tab:analogues}.

\noindent\scriptsize
\setlength{\tabcolsep}{3pt}
\renewcommand{\arraystretch}{1.18}
\begin{longtable}{|>{\RaggedRight\arraybackslash}p{0.14\linewidth}|>{\RaggedRight\arraybackslash}p{0.20\linewidth}|>{\RaggedRight\arraybackslash}p{0.28\linewidth}|>{\RaggedRight\arraybackslash}p{0.28\linewidth}|}
\caption{Сравнение существующих шахматных онлайн-платформ}\label{tab:analogues}\\
\hline
Решение & Назначение & Характерные функции & Что учтено в ChessView\\ \hline
\endfirsthead
\hline
Решение & Назначение & Характерные функции & Что учтено в ChessView\\ \hline
\endhead
Chess.com & Массовая коммерческая шахматная платформа & Онлайн-игра, рейтинги, задачи, обучение, турниры, профиль, социальные функции & Связка профиля, live-партии, рейтинга, истории и задач\\ \hline
Lichess & Бесплатная шахматная платформа с открытой идеологией & Быстрые партии, анализ, задачи, турниры, публичные API & Разделение live-игры и анализа, удобная работа с FEN/PGN\\ \hline
Tornelo & Турнирная платформа & Регистрация участников, раунды, пары, расписания, OTB-сценарии & Турнирный модуль, Swiss helpers, OTB manager и scheduled matches\\ \hline
ChessView & Учебное full-stack веб-приложение & Auth, live game, WebSocket, анализ, задачи, турниры, админка, база данных & Демонстрация архитектуры, БД, API, клиента и админки без заявления о промышленном масштабе\\ \hline
\end{longtable}
\normalsize

Анализ аналогов показывает, что для дипломного проекта важно не копировать весь объем крупных платформ, а выбрать ядро предметной области. В ChessView таким ядром являются авторизация, профиль, партия, ход, время, история, задача и турнир. Дополнительные поверхности, такие как магазин, клубы и OTB-менеджер, используются как расширение пользовательского интерфейса и демонстрация роста платформы.

\section{Инструменты разработки}

Инструменты разработки выбраны так, чтобы каждая часть проекта решала свою задачу. Backend написан на Python и FastAPI, frontend реализован на React и TypeScript, база данных работает на PostgreSQL, миграции управляются Alembic, контейнерный запуск описан в Docker Compose. Такой набор инструментов распространен в современной веб-разработке и хорошо подходит для учебного full-stack проекта.

Python --- интерпретируемый высокоуровневый язык программирования с динамической типизацией и развитой стандартной библиотекой. В проекте он используется на серверной стороне. Python удобен для быстрого описания бизнес-логики, интеграции с библиотеками и написания служебных скриптов. В ChessView важна также поддержка асинхронного программирования: FastAPI, SQLAlchemy async и WebSocket-обработчики используют \code{async}/\code{await}, что позволяет серверу обрабатывать сетевые операции без блокирования процесса. Официальная документация Python приведена в источнике [1].

FastAPI --- веб-фреймворк для Python, ориентированный на создание API. Он использует типы Python и Pydantic-схемы, автоматически формирует OpenAPI-описание, поддерживает dependency injection и асинхронные обработчики. В ChessView FastAPI регистрирует REST routers, WebSocket endpoint, CORS, обработчики исключений и healthcheck. Документация FastAPI указана в источнике [2].

Uvicorn --- ASGI-сервер, который запускает FastAPI-приложение. Он принимает HTTP и WebSocket-соединения и передает их приложению. В режиме разработки Uvicorn запускается с параметром \code{--reload}, чтобы сервер перезапускался после изменения кода. Описание Uvicorn и ASGI приведено в источнике [3].

PostgreSQL --- реляционная система управления базами данных. В проекте она хранит пользователей, партии, ходы, турниры, задачи, платежные события и административный аудит. PostgreSQL выбран из-за поддержки транзакций, внешних ключей, индексов, UUID, JSON-полей и надежной работы с отношениями между таблицами. Документация PostgreSQL приведена в источнике [4].

SQLAlchemy --- Python-библиотека для работы с базами данных и объектно-реляционного отображения. В ChessView SQLAlchemy используется в асинхронном режиме: инфраструктурные репозитории преобразуют доменные сущности в ORM-модели и обратно. Это позволяет отделить бизнес-логику от конкретных SQL-запросов. Документация SQLAlchemy указана в источнике [5].

Alembic --- инструмент миграций для SQLAlchemy. Он хранит историю изменения схемы: baseline, добавление внешних ключей, расширение платформы, платежи, wallet, speed ratings, passkeys и face templates. Благодаря Alembic схема базы данных воспроизводима и соответствует версии backend-кода. Документация Alembic приведена в источнике [6].

TypeScript --- язык программирования, расширяющий JavaScript статической типизацией. В клиентской части TypeScript помогает описывать структуры ответов API, состояние пользователя, игровые сущности и props компонентов. Это снижает количество ошибок при изменении контракта между frontend и backend. Документация TypeScript приведена в источнике [7].

React --- библиотека для создания пользовательских интерфейсов. В ChessView React используется для построения SPA: страниц авторизации, лобби, партии, анализа, профиля, турниров, задач и административной панели. Компонентный подход позволяет выделять страницы, виджеты, features и общие UI-примитивы. Документация React указана в источнике [8].

Vite --- инструмент сборки frontend-проектов. Он обеспечивает быстрый dev server, сборку production-бандла и proxy для запросов \code{/api}, \code{/media} и \code{/ws}. Благодаря этому frontend может работать на \code{localhost:5173}, а запросы передаются backend на \code{localhost:8000}. Документация Vite приведена в источнике [9].

TanStack Query используется для загрузки и кэширования серверных данных на frontend. Zustand используется для легкого клиентского состояния, например для авторизации. React Router описывает маршруты и protected routes. chess.js и react-chessboard помогают отображать шахматные позиции на клиенте, а Stockfish используется для локального анализа. Эти инструменты указаны в источниках [10]--[14].

Docker Compose используется для воспроизводимого запуска нескольких сервисов: PostgreSQL, backend, frontend и admin-frontend. Он фиксирует порты, зависимости контейнеров, переменные окружения и volume для базы данных. Такой запуск удобен для демонстрации дипломного проекта, потому что уменьшает количество ручных действий. Документация Docker Compose приведена в источнике [15].

\noindent\scriptsize
\setlength{\tabcolsep}{3pt}
\renewcommand{\arraystretch}{1.18}
\begin{longtable}{|>{\RaggedRight\arraybackslash}p{0.20\linewidth}|>{\RaggedRight\arraybackslash}p{0.28\linewidth}|>{\RaggedRight\arraybackslash}p{0.36\linewidth}|>{\RaggedRight\arraybackslash}p{0.10\linewidth}|}
\caption{Инструменты разработки и их назначение}\label{tab:tools}\\
\hline
Инструмент & Где используется & Конкретная роль в ChessView & Источник\\ \hline
\endfirsthead
\hline
Инструмент & Где используется & Конкретная роль в ChessView & Источник\\ \hline
\endhead
Python 3.12 & Backend & Асинхронная серверная логика, доменные сервисы и служебные скрипты & [1]\\ \hline
FastAPI & Backend & REST API, WebSocket endpoint, DI, схемы ответов & [2]\\ \hline
Uvicorn & Backend runtime & Запуск ASGI-приложения & [3]\\ \hline
PostgreSQL 16 & Database & Реляционное хранение данных & [4]\\ \hline
SQLAlchemy async & Backend infrastructure & ORM-модели и репозитории & [5]\\ \hline
Alembic & Database lifecycle & Миграции схемы БД & [6]\\ \hline
TypeScript & Frontend/admin & Типизация UI и API-контрактов & [7]\\ \hline
React & Frontend/admin & Компонентный интерфейс приложения & [8]\\ \hline
Vite & Frontend tooling & Dev server, proxy, production build & [9]\\ \hline
TanStack Query & Frontend data & Кэширование и обновление серверных данных & [10]\\ \hline
Zustand & Frontend state & Локальное состояние авторизации & [11]\\ \hline
React Router & Frontend routing & Маршруты и защищенные страницы & [12]\\ \hline
python-chess & Backend game rules & Проверка шахматных ходов & [13]\\ \hline
Stockfish & Browser analysis & Локальная оценка позиций & [14]\\ \hline
Docker Compose & Deployment & Запуск Postgres, backend, frontend и admin & [15]\\ \hline
ESLint & Frontend tooling & Статическая проверка TypeScript и React-кода & [16]\\ \hline
\end{longtable}
\normalsize



Выбор инструментов связан с фактическими задачами проекта, а не с формальным перечислением популярных технологий. Python и FastAPI выбраны для backend, потому что сервер должен работать с HTTP, WebSocket, асинхронной базой данных и доменными сервисами. PostgreSQL выбран для БД, потому что проект хранит связанные сущности и нуждается во внешних ключах. TypeScript и React выбраны для frontend, потому что интерфейс содержит много состояний: авторизация, загрузка профиля, игровая доска, очередь, анализ, турниры и admin-сессия.

Средства разработки также разделены по жизненному циклу. На этапе написания кода используются TypeScript, Python, React, FastAPI и SQLAlchemy. На этапе хранения структуры БД используется Alembic. На этапе запуска используется Uvicorn или Docker Compose. На этапе frontend-сборки используется Vite. На этапе обращения к источникам используются официальные документации инструментов, поэтому ссылки на них вынесены в список источников и указаны в тексте раздела.

Важно, что инструменты не дублируют друг друга. SQLAlchemy не заменяет PostgreSQL, а предоставляет программный слой для работы с ней. Alembic не хранит данные приложения, а хранит историю изменения схемы. React не выполняет серверную проверку шахматных правил, а показывает интерфейс и отправляет действия. Stockfish не используется для проверки ходов live-партии, а применяется в учебном анализе позиции. Такое разграничение делает стек понятнее и уменьшает количество спорных мест в архитектуре.

\section{Требования к серверу для запуска приложения}

ChessView рассчитан на локальный демонстрационный запуск и single-instance deployment. Это означает, что backend, WebSocket rooms, matchmaking queue и игровой монитор работают в рамках одного процесса FastAPI. Для промышленного горизонтального масштабирования потребовались бы Redis или аналогичное shared-хранилище для очередей и WebSocket fanout, распределенная координация фоновых задач и объектное хранилище для media-файлов.

Минимальные требования к серверу приведены в таблице~\ref{tab:server-requirements}. Они учитывают запуск через Docker Compose и одновременную работу PostgreSQL, backend, frontend и административного frontend.

\noindent\scriptsize
\setlength{\tabcolsep}{4pt}
\renewcommand{\arraystretch}{1.18}
\begin{longtable}{|>{\RaggedRight\arraybackslash}p{0.22\linewidth}|>{\RaggedRight\arraybackslash}p{0.30\linewidth}|>{\RaggedRight\arraybackslash}p{0.40\linewidth}|}
\caption{Требования к серверу для запуска ChessView}\label{tab:server-requirements}\\
\hline
Параметр & Минимальное значение & Пояснение\\ \hline
\endfirsthead
\hline
Параметр & Минимальное значение & Пояснение\\ \hline
\endhead
ОС & Windows 10/11, Linux или macOS & Нужна поддержка Docker или локального запуска Python/Node/PostgreSQL\\ \hline
CPU & 2 ядра & Достаточно для демонстрации backend, frontend и БД\\ \hline
RAM & 4 ГБ & Для комфортной работы Docker Compose лучше 8 ГБ\\ \hline
Диск & 2 ГБ свободного места & Код, зависимости, Docker layers, база и media-файлы\\ \hline
Порты & 5432, 8000, 5173, 5174 & PostgreSQL, backend API, пользовательский frontend, admin frontend\\ \hline
Docker & Docker Engine + Compose & Основной воспроизводимый запуск проекта\\ \hline
Локальный backend & Python 3.12, uv & Альтернативный split-запуск backend\\ \hline
Локальный frontend & Node.js, Yarn Classic & Альтернативный split-запуск frontend и admin frontend\\ \hline
Переменные окружения & \code{.env} & DATABASE\_URL, JWT\_SECRET, CORS, VITE\_SERVER\_URL и proxy targets\\ \hline
\end{longtable}
\normalsize

Для демонстрации используется команда \code{docker compose up --build}. После запуска пользовательский интерфейс доступен на \code{http://localhost:5173}, административный интерфейс на \code{http://localhost:5174}, backend healthcheck на \code{http://localhost:8000/health}, а API на \code{http://localhost:8000/api/v1}. Такая схема запуска показана на рисунке~\ref{fig:deploy-topology}.



Серверный запуск ChessView можно рассматривать в двух режимах. Первый режим — демонстрационный, через Docker Compose. Он предпочтителен для защиты, потому что одной командой поднимает все сервисы и уменьшает риск забыть отдельный процесс. Второй режим — split development. В нем PostgreSQL может запускаться в контейнере, backend запускается через uv и Uvicorn, frontend через Yarn, а admin-frontend отдельной командой. Этот режим удобен при разработке и отладке.

Для запуска важно заранее проверить файл .env. В нем должны быть заданы учетные данные PostgreSQL, DATABASE_URL для локального backend, DOCKER_DATABASE_URL для backend внутри контейнера, JWT_SECRET, параметры времени жизни токенов, CORS_ORIGINS, BACKEND_URL, FRONTEND_URL и VITE_SERVER_URL. Если эти значения не согласованы, frontend может открыться, но API-запросы будут уходить не туда или backend не подключится к базе данных.

Требование single-instance также должно быть явно зафиксировано. Текущая версия ChessView рассчитана на один backend-процесс. Очередь matchmaking, WebSocket rooms, connection manager и background game monitor находятся в памяти процесса. Это допустимо для дипломного проекта и локальной демонстрации. Для горизонтального масштабирования понадобились бы Redis или аналогичная shared-система, распределенная блокировка фоновых задач и общий объектный storage для media-файлов.

\begin{figure}[H]
\centering
\begin{adjustbox}{max width=0.95\textwidth,max height=0.30\textheight}
\begin{tikzpicture}[node distance=13mm, box/.style={draw, rounded corners=2pt, align=center, minimum width=34mm, minimum height=9mm, fill=gray!8}, arr/.style={-{Stealth[length=2mm]}, thick}]
\node[box] (browser) {Браузер\\пользователя};
\node[box, right=18mm of browser] (front) {frontend\\localhost:5173};
\node[box, below=10mm of front] (admin) {admin-frontend\\localhost:5174};
\node[box, right=18mm of front] (back) {backend FastAPI\\localhost:8000};
\node[box, right=18mm of back] (db) {PostgreSQL\\localhost:5432};
\draw[arr] (browser) -- (front);
\draw[arr] (browser) -- (admin);
\draw[arr] (front) -- node[above, font=\scriptsize]{/api, /ws} (back);
\draw[arr] (admin) -- node[below, font=\scriptsize]{/api/v1/admin} (back);
\draw[arr] (back) -- (db);
\end{tikzpicture}
\end{adjustbox}
\caption{Топология локального запуска ChessView}\label{fig:deploy-topology}
\end{figure}
